% !TeX root = ../papers/nc-p3.tex
\section{可動平面束の排除と全曲面射影の被覆}
\label{sec:exclude-moving-plane-pencils}
\begin{theorem}\label{thm:no-moving-plane-pencils}
標準正則窓では各辺の両射影で直線ファイバーの中心は有限個である。
従って一辺でも射影像が二次元以上となる全窓は定理\ref{thm:isolated-surface-fibers}の分類で尽きる。
\end{theorem}
\begin{proof}[二つの平面束]
可動中心の直線 $L$ を位置0に取り、位置1の後続平面の方程式を $b$ とする。
系\ref{cor:rank-three-plane-branch-slice}の表示を使い、反対辺 $R_2$ の純関係を $c\otimes d$ とする。
$R_1$ は $c$ に終わる純関係を持たない。従って $V_2/kc$ へ射影した六関係は独立で、
さらに $V_1/kb$ へ射影しても、失う次元は高々3である。
一方 $w$ を $L$ に制限した後の $R_1$ 側の縮約像は、表示式より
$b\otimes V_2+ V_1\otimes c$ に入るので全て失われる。
ゆえにこの縮約像、すなわち逆辺 $R_3$ の $L$ への制限の階数は高々3である。
ここには $d\otimes H^0(\mathcal O_L(1))$ が入る（独立な二つの動く係数を縮約する）。
各 $x\in L$ で評価した空間は次元2以上であり、従ってちょうど2である。
逆辺にも動く直線束があり、定理\ref{thm:all-endpoint-ranks-two}による支持の一意性と
平面束の分類から、その平面方程式は共通の文字 $d$ である。

$x=(x_0,x_1)$ として順方向の支持を $\langle b,l(x)\rangle$、逆方向を
$\langle d,m(x)\rangle$ とする。$l,m$ の二係数はそれぞれ平面方程式を法として独立である。
厳密四周期表示で $w$ の $L$ への制限は、位置1,2,3の順に
\[
 b\otimes A(x)\otimes d+
 l(x)\otimes c\otimes d+
 b\otimes e\otimes m(x)
\]
と書ける。実際、各支持束は定数直線と $\mathcal O(-1)$ の直和で、
二つの動く因子を同時に使う項は $\mathcal O(-2)(1)$ の切断なので零である。
$c,e$ は非零である（各点の両端階数は2）。$R_1$ は $b\otimes e$ を含むため
$R_0$ に $b$ に終わる純関係はなく、$R_2$ に $e$ に始まる純関係もない。

$R_2$ で $d$ に終わる関係は正確に一本である。実際 $V_2/ke$ を先に取り次に $V_3/kd$ を取ると、
初めは六次元で次も三次元以上である。逆に $d$ を先に消すと、今の表示を位置0,1で縮約して
得る $e\otimes m(x)$ の二係数が射影像にあり、次に $e$ を消すと少なくとも二次元を失う。
従って $d$ だけを消した像は五次元以上で、$c\otimes d$ を失うから正確に五次元である。
\end{proof}

\begin{proof}[二つの独立な束パラメータを持つ三文字]
位置0の相続く二回の訪問で $x,x'\in L$ を別々に指定する。
その間の位置1と3の残る文字はそれぞれ
\[
 E_x=V_1/\langle b,l(x)\rangle,\qquad F_{x'}=V_3/\langle d,m(x')\rangle .
\]
束として $E,F\simeq\mathcal O\oplus\mathcal O(1)$ である。
$R_1$ の $E\otimes V_2$ への像束を $I$、$R_2$ の $V_2\otimes F$ への像束を $J$ とする。
系\ref{cor:endpoint-line-line-cut}により両方の核は各点で二次元であり、
表示式から第一の核は $b\otimes e$ と $b\otimes A(x)+l(x)\otimes c$ で張られる。
第二の核も同様に定数一本と一次の一本である。よって $I,J$ は階数4、次数1である。
$L\times L$ 上で記号の左右の束を各射影で引き戻す。挿入写像
\[
 D:(I\boxtimes F)\oplus(E\boxtimes J)
 \longrightarrow E\boxtimes F\otimes V_2
\]
では三文字を元の順に並べる。両側の階数は16、行列式は
$\mathcal O(2,2)$ の切断である。例えば第一因子での次数は標的8、始域 $2+4=6$ である。
余核は端点を指定した三文字の商であり、対角上の各点でその次元は高々1である
（同じ点を四文字ごとに繰り返したときの独立ブロックの増大による）。

対角上の一般点で余核は非零である。反対に一般点で全射なら、そこでは二つの挿入の交叉がないため、
全ての三次重なりは両端を $E_x,F_x$ に落とすと零になる。
$Y=V_1/kb,\ U=V_3/kd$ と書く。$Y\otimes U\to E_x\otimes F_x$ の全 $x$ に共通の核は
$l,m$ のそれぞれ二係数の張る二空間どうしに入る。実際、この二空間の補直線を取り、
補直線どうしの成分と片側だけ補直線の成分は商の対応する直和成分を全 $x$ で評価すると検出できる。
従って $w$ を $Y,U$ に射影したものの $Y$ 支持は高々二次元である。
さらにその位置0支持は $L^\perp$ に入る（先の三項の表示は $b,d$ を消せば零）。
すると位置0,1対位置2,3の階数は高々4である。
しかし先に $b$ を消すと $R_0$ 側に核はなく階数6であり、続いて $d$ を消すと
$R_2$ 側に失うのは一本だけなので階数5である。これは矛盾する。
従って対角全点で $D$ の階数は15である。
\end{proof}

\begin{proof}[二重対応に沿った指数増大]
行列式が恒等零なら $L$ の異なる二点の任意の順序対で三文字商が非零である。
各対ごとにその非零な一次元商を取り、四文字ごとにこの二点を任意に並べる。
各境界で次に選んだ二点は一次式で分離できるので、長さ $4n+1$ で少なくとも $2^n$ の次元を得る
（以下の経路評価の被約二点版でもある）。

行列式が恒等零でなければ、零因子は対角と型 $(1,1)$ の残余因子の和である。
第一射影の垂直成分を取り除いた有効因子を $C\subset L\times L$ とする。
第一射影は有限平坦次数2である。各ファイバーは射影直線上の非零な二次式の零スキームだからである。
$C$ は二重対角、対角と自己射影同型のグラフ、または対角と定数グラフの和である。
最後の二型では共通成分がないものとする。行列式の単純成分の一般点で $D$ は階数15である。
階数14以下なら行列式はその成分の局所方程式の二乗で割れるからである。
また二重対角を含め、対角上では既に階数15である。

各 $n$ で一般初期点 $x_0$ を一つ固定し、$(x_{j-1},x_j)\in C$ となる
長さ $n$ のスキームを $T_n$ とする。有限平坦性を逐次引き戻すと長さは $2^n$ である。
全ての経路の全辺で $D$ は階数15の開集合に入るようにできる。
自己同型の場合には有限回の合成で有限例外集合を避ければよく、
定数グラフで一度定数点に入った後はその対角点しか通らない。
その開集合上では行列式が唯一の残る関係で、$D$ の余核の $C$ への引戻しは可逆層である。
これは可逆な15次小行列を使う行・列消去から、非被約の場合にも成立する。

各辺で三文字をこの可逆商へ評価し、境界の位置0の文字を $x_j$ で $\mathcal O(1)$ へ評価する。
得られる積評価は元の全二次関係を消す。内部は $I,J$ を割ってあり、
境界の $R_0,R_3$ はそれぞれ $E_{x_j},F_{x_j}$ の定義で零である。
各内部三文字の線型結合を、有限個の支持点全部で非零になるように選べる
（余核は三文字で生成され、体は無限）。これらと初期点の非零評価を固定すると可逆因子を掛けるだけとなる。
残る $n$ 個の境界の二一次式の積は $T_n$ 上の対応する可逆層の全切断を張る。
実際、各 $T_j\to T_{j-1}$ の各ファイバーは射影直線の長さ2の部分スキームで、
二一次式の制限はその $\mathcal O(1)$ の切断への同型である。
有限平坦性と中山の補題により直像束はこの二切断を基底に持ち、射影公式で帰納できる。
従って長さ $4n+1$ の経路空間の次元は少なくとも $2^n$ である。
この下界は二重点でも二つの異なる幾何学的後続を要しない。標準Hilbert関数の三次増大と矛盾する。
これで可動直線束は全て排除された。
\end{proof}
